\chapter{Introducción y Objetivos}
\label{chap:introduccion}

En el presente trabajo práctico, se determinaron las características
fundamentales de una red de dos puertos mediante el empleo de métodos de
medición e instrumental de laboratorio convencionales. Específicamente, se
trabajó sobre un amplificador de baja frecuencia y un amplificador de dos etapas
con realimentación negativa, caracterizando en ambos sus parámetros básicos,
tales como las impedancias de entrada y salida, la potencia de salida, y las
ganancias tanto de tensión como de potencia.

Como parte del estudio de estos parámetros, se evaluó el impacto de la
realimentación negativa sobre la ganancia y las impedancias del sistema,
contrastando experimentalmente su comportamiento en condiciones de lazo abierto
y lazo cerrado. Esta etapa permitió evidenciar las ventajas y desventajas
cualitativas de la realimentación, así como validar la aplicabilidad de los
métodos de medición ensayados.

Para la determinación de la impedancia de entrada ($Z_i$) y de salida ($Z_o$) de
los amplificadores se implementó el método clásico de reducción de la tensión a
la mitad mediante una carga variable ($R_L = R_o$) y se analizaron críticamente
las limitaciones de este enfoque en sistemas realimentados; específicamente, se
evaluó cómo la tendencia del lazo cerrado a la estabilización de la tensión de
salida altera la dinámica de la medición, considerando asimismo los riesgos
asociados a la sobreexcitación y la disipación de potencia en los componentes.

Asimismo, se analizó la respuesta en frecuencia de los sistemas mediante dos
enfoques metodológicos complementarios. Por un lado, se efectuó un barrido en
frecuencia para obtener la curva característica en régimen permanente y, por el
otro, se evaluó la respuesta transitoria mediante la aplicación de una onda
cuadrada y la posterior medición de su tiempo de crecimiento ($t_r$).

A lo largo de todo el desarrollo experimental, se incorporaron y aplicaron
conceptos fundamentales de la práctica metrológica, tales como el uso y la
utilidad de las escalas en decibelios ($dB$) y la implementación del método de
sustitución para la realización de mediciones indirectas.

Por lo tanto, los objetivos del trabajo práctico son:
\begin{itemize}
    \item Determinar los parámetros fundamentales de un amplificador: impedancia de entrada ($Z_i$),
          impedancia de salida ($Z_o$), ganancia de tensión ($A_v$) y ganancia de potencia ($A_p$).
    \item Adquirir destreza en el uso de instrumentos de medición con escalas graduadas en decibeles ($dB$),
          comprendiendo el significado físico de las unidades $dBu$ y $dBm$.
    \item Emplear el osciloscopio como herramienta principal para efectuar el análisis dinámico y la
          medición de parámetros temporales y de amplitud en sistemas de amplificación.
    \item Evaluar experimentalmente el efecto de la realimentación negativa sobre la ganancia, la respuesta
          en frecuencia (ancho de banda), la distorsión y las impedancias de entrada y salida.
    \item Expresar los resultados obtenidos en las mediciones acompañados de su correspondiente grado de
          incertidumbre experimental.
\end{itemize}


